\section{Experiments}
\label{sec:experiments}

To evaluate the Set Transformer, we apply it to a suite of tasks involving sets of data points.
%We applied the Set Transformer to various problems that use sets of data.
We repeat all experiments five times and report performance metrics evaluated on corresponding test datasets.
Along with baselines, we compared various architectures arising from the combination of the choices of having attention in encoders and decoders. %, each of which roughly represents existing works as its special cases.
Unless specified otherwise, ``simple pooling" means average pooling.
\begin{itemize}
\setlength\itemsep{-0.4em}
\item rFF + Pooling \citep{Zaheer2017}: rFF layers in encoder and simple pooling + rFF layers in decoder.
\item rFFp-mean/rFFp-max + Pooling \citep{Zaheer2017}: rFF layers with permutation equivariant variants in encoder \citep[\eqref{eq:multihead1}]{Zaheer2017} and simple pooling + rFF layers in decoder.
\item rFF + Dotprod \citep{Yang2018,Ilse2018}: rFF layers in encoder and dot product attention based weighted sum pooling + rFF layers in decoder.
\item SAB (ISAB) + Pooling (ours):  Stack of SABs (ISABs) in encoder and simple pooling + rFF layers in decoder.
\item rFF + PMA (ours): rFF layers in encoder and PMA (followed by stack of SABs) in decoder.
\item SAB (ISAB) + PMA (ours): Stack of SABs (ISABs) in encoder and PMA (followed by stack of SABs) in decoder.
\end{itemize}

\subsection{Toy Problem: Maximum Value Regression}
\begin{table}[t]
\centering
\small
\caption{Mean absolute errors on the max regression task.}
\vspace{5pt}
\begin{tabular}{@{}ccccc@{}}
\toprule
Architecture & MAE \\\midrule
rFF + Pooling ($\mathrm{mean}$) & 2.133 $\pm$ 0.190 \\
rFF + Pooling ($\mathrm{sum}$) & 1.902 $\pm$ 0.137 \\
rFF + Pooling ($\mathrm{max}$) & \bf 0.1355 $\pm$ 0.0074\\
\midrule
SAB + PMA (ours) & 0.2085 $\pm$ 0.0127\\
\bottomrule
\label{table:max}
\end{tabular}
\end{table}
\label{subsec:max_regression}
To demonstrate the advantage of attention-based set aggregation over simple pooling operations, we consider a toy problem: regression to the maximum value of a given set.
Given a set of real numbers $\{ x_1, \ldots, x_n\}$, the goal is to return $\mathrm{max}(x_1, \cdots, x_n)$.
Given prediction $p$, we use the mean absolute error $|p - \mathrm{max}(x_1, \cdots, x_n)|$ as the loss function. We constructed simple pooling architectures with three different pooling operations: $\operatorname{max}$, $\operatorname{mea}$n, and $\operatorname{sum}$. We report loss values after training in Table~\ref{table:max}.
%Note that even though the max pooling architecture only has to learn the identity function to predict the output perfectly, Set Transformers achieve comparable performance.
%This is likely because Set Transformers can learn to find and attend to the maximum element.
%Note that the mean- and sum-pooling architectures have far higher mean absolute errors (MAE), which indicates
Mean- and sum-pooling architectures result in a high mean absolute error (MAE).
%The model with max-pooling only has to learn the identity function as its encoder to predict the output perfectly, and thus achieves the highest performance.
The model with max-pooling can predict the output perfectly by learning its encoder to be an identity function, and thus achieves the highest performance.
Notably, the Set Transformer achieves performance comparable to the max-pooling model, which underlines the importance of additional flexibility granted by attention mechanisms --- it can learn to find and attend to the maximum element.

%Our analysis of this is that learning interactions among items is difficult for aggregation functions such as sum and mean.
%The max regression problem requires knowledge of such interactions because the maximum value of a set $\{ x_1, \cdots, x_n\}$ only depends on one maximum element and thus the network must at least learn how to compare items.
%We would like to point out that most pooling architectures for sets use mean-pooling.


\begin{figure}[t]
\centering\includegraphics[width=\columnwidth]{figs/omni20.png}
\caption{Counting unique characters: this is a randomly sampled set of 20 images from the Omniglot dataset. There are 14 different characters inside this set.}
\label{fig:omni_try}
\end{figure}

\begin{table}[t]
\centering
\small
\caption{Accuracy on the unique character counting task.}
\vspace{5pt}
\begin{tabular}{@{}cc@{}} \toprule
Architecture & Accuracy \\ \midrule
rFF + Pooling & 0.4382 $\pm$ 0.0072 \\
rFFp-mean + Pooling & 0.4617 $\pm$ 0.0076 \\
rFFp-max + Pooling & 0.4359 $\pm$ 0.0077 \\
rFF + Dotprod & 0.4471 $\pm$ 0.0076 \\
\midrule
rFF + PMA (ours) & 0.4572 $\pm$ 0.0076\\
SAB + Pooling (ours) & 0.5659 $\pm$ 0.0077 \\
SAB + PMA (ours) & \bf 0.6037 $\pm$ 0.0075 \\
\bottomrule
\label{table:unique}
\end{tabular}
\end{table}
\begin{figure}[t]
\centering
\includegraphics[width=0.90\linewidth]{figs/isab_n.pdf}
\caption{
Accuracy of $\mathrm{ISAB}_n + \mathrm{PMA}$ on the unique character counting task.
x-axis is $n$ and y-axis is accuracy.
}
\label{fig:isab_n}
\end{figure}

\begin{table*}[h]
	\centering
	\small
	\caption{Meta clustering results.
	The number inside parenthesis indicates the number of inducing points used in ISABs of encoders.
	We show average likelihood per data for the synthetic dataset and the adjusted rand index (ARI) for the CIFAR-100 experiment.
	LL1/data, ARI1 are the evaluation metrics after a single EM update step.
	The oracle for the synthetic dataset is the log likelihood of the actual parameters used to generate the set, and the CIFAR oracle was computed by running EM until convergence.
	%The results measured from function outputs (LL0/data, ARI0) and the results after single EM update (LL1/data, ARI1) are presented.
	}
	\vspace{5pt}
	\begin{tabular}{@{}ccccc@{}}\toprule
		& \multicolumn{2}{c}{Synthetic} & \multicolumn{2}{c}{CIFAR-100} \\
		\cmidrule{2-3}
		\cmidrule{4-5}
	Architecture & LL0/data & LL1/data & ARI0 & ARI1  \\
	\midrule
	Oracle & -1.4726 & & 0.9150 &\\
	rFF + Pooling & -2.0006 $\pm$ 0.0123 &-1.6186 $\pm$ 0.0042 & 0.5593 $\pm$ 0.0149 & 0.5693 $\pm$ 0.0171 \\
	rFFp-mean + Pooling & -1.7606 $\pm$ 0.0213 &-1.5191 $\pm$ 0.0026 & 0.5673 $\pm$ 0.0053 &0.5798 $\pm$ 0.0058 \\
	rFFp-max + Pooling & -1.7692 $\pm$ 0.0130 &-1.5103 $\pm$ 0.0035 & 0.5369 $\pm$ 0.0154 &0.5536 $\pm$ 0.0186 \\
	rFF + Dotprod & -1.8549 $\pm$ 0.0128 &-1.5621 $\pm$ 0.0046 & 0.5666 $\pm$ 0.0221 &0.5763 $\pm$ 0.0212 \\
    \midrule
	SAB + Pooling (ours) & -1.6772 $\pm$ 0.0066 &-1.5070 $\pm$ 0.0115 & 0.5831 $\pm$ 0.0341 &0.5943 $\pm$ 0.0337\\
	ISAB (16) + Pooling (ours) & -1.6955 $\pm$ 0.0730 &-1.4742 $\pm$ 0.0158 & 0.5672 $\pm$ 0.0124 &0.5805 $\pm$ 0.0122\\
	rFF + PMA (ours) & -1.6680 $\pm$ 0.0040 &-1.5409 $\pm$ 0.0037 &  0.7612 $\pm$ 0.0237 &0.7670 $\pm$ 0.0231\\
	SAB + PMA (ours) & -1.5145 $\pm$ 0.0046 &-1.4619 $\pm$ 0.0048 &  0.9015 $\pm$ 0.0097 &0.9024 $\pm$ 0.0097\\
	ISAB (16) + PMA (ours) & \bf -1.5009 $\pm$ 0.0068 & \bf -1.4530 $\pm$ 0.0037 & \bf 0.9210 $\pm$ 0.0055 &\bf 0.9223 $\pm$ 0.0056 \\
	\bottomrule
	\end{tabular}
	\label{tab:meta_clustering}
\end{table*}

\begin{figure*}[t]
	\centering
	\includegraphics[width=0.4\linewidth]{figs/ff.pdf}
	\includegraphics[width=0.4\linewidth]{figs/sab_ff.pdf}
	\includegraphics[width=0.4\linewidth]{figs/ff_sab.pdf}
	\includegraphics[width=0.4\linewidth]{figs/sab.pdf}
	\caption{
    Clustering results for 10 test datasets, along with centers and covariance matrices.
    rFF+Pooling (top-left), SAB+Pooling (top-right), rFF+PMA (bottom-left), Set Transformer (bottom-right).
    Best viewed magnified in color.}
	\label{fig:synthetic_clustering}
\end{figure*}

\subsection{Counting Unique Characters}
\label{subsec:unique}
In order to test the ability of modelling interactions between objects in a set, we introduce a new task of counting unique elements in an input set.
We use the Omniglot~\citep{Lake2015} dataset, which consists of 1,623 different handwritten characters from various alphabets, where each character is represented by 20 different images.

We split all characters (and corresponding images) into train, validation, and test sets and only train using images from the train character classes.
We generate input sets by sampling between 6 and 10 images and we train the model  to predict the number of different characters inside the set.
We used a Poisson regression model to predict this number, with the rate $\lambda$ given as the output of a neural network.
We maximized the log likelihood of this model using stochastic gradient ascent.

We evaluated model performance using sets of images sampled from the test set of characters.
%Therefore, not only has the network not seen any of the images inside a given test input, but it has not even seen any image of the same character class during training.
Table~\ref{table:unique} reports accuracy, measured as the frequency at which the mode of the Poisson distribution chosen by the network is equal to the number of characters inside the input set.

We additionally performed experiments to see how the number of incuding points affects performance.
We trained $\mathrm{ISAB}_n + \mathrm{PMA}$ on this task while varying the number of inducing points ($n$).
Accuracies are shown in Figure \ref{fig:isab_n}, where other architectures are shown as horizontal lines for comparison.
Note first that even the accuracy of $\mathrm{ISAB}_1 + \mathrm{PMA}$ surpasses that of both $\mathrm{rFF} + \mathrm{Pooling}$ and $\mathrm{rFF} + \mathrm{PMA}$, and that performance tends to increase as we increase $n$.

\subsection{Amortized Clustering with Mixture of Gaussians}
%We ran all the algorithms five times and measured the averages.
We applied the set-input networks to the task of maximum likelihood of mixture of Gaussians (MoGs).
The log-likelihood of a dataset $X = \{x_1, \dots, x_n\}$ generated from an MoG with $k$ components is
{
\[
\log p(X; \theta)
= \sum_{i=1}^n  \log \sum_{j=1}^k \pi_j \gN ( x_i ; \mu_j, \mathrm{diag}(\sigma^2_j)).
\]}
The goal is to learn the optimal parameters $\theta^*(X) = \argmax_\theta \log p(X;\theta)$.
The typical approach to this problem is to run an iterative algorithm such as Expectation-Maximisation (EM) until convergence.
Instead, we aim to learn a generic meta-algorithm that directly maps the input set $X$ to $\theta^*(X)$.
One can also view this as amortized maximum likelihood learning.
Specifically, given a dataset $X$, we train a neural network to output parameters $f(X;\lambda) = \{\pi(X), \{\mu_j(X), \sigma_j(X)\}_{j=1}^k\}$ which maximize
{\small
\[
\E_X \left[ \sum_{i=1}^{|X|} \log \sum_{j=1}^k \pi_j(X) \gN(x_i ; \mu_j(X), \mathrm{diag}(\sigma_j^2(X)))\right].
\]}
We structured $f(\cdot;\lambda)$ as a set-input neural network and learned its parameters $\lambda$ using stochastic gradient ascent, where we approximate gradients using minibatches of \emph{datasets}.

We tested Set Transformers along with other set-input networks on two datasets.
We used four seed vectors for the PMA ($S\in\R^{4\times d}$) so that each seed vector generates the parameters of a cluster.

\textbf{Synthetic 2D mixtures of Gaussians}:  Each dataset contains $n \in [100, 500]$ points on a 2D plane, each sampled from one of four Gaussians.
%Average log-likelihhod/data over 1,000 test sets were compared.
%The oracle was computed by measuring log-likelihood w.r.t. the true parameters used to generate test sets.

\textbf{CIFAR-100}: Each dataset contains $n \in [100, 500]$ images sampled from four random classes in the CIFAR-100 dataset.
Each image is represented by a 512-dim vector obtained from a pretrained VGG network~\citep{Simonyan2014}.
%Adjusted Rand Index (ARI) over 1,000 test sets were compared, and the oracle was computed by running EM for each test set independently.

\begin{table*}[t]
\centering
\small
\caption{
Test accuracy for the point cloud classification task using $100, 1000, 5000$ points.
}
\vspace{5pt}
\begin{tabular}{@{}ccccc@{}}
\toprule
Architecture & 100 pts & 1000 pts & 5000 pts \\
\midrule
rFF + Pooling \citep{Zaheer2017} & - & 0.83 $\pm$ 0.01 & - \\
rFFp-max + Pooling \citep{Zaheer2017} & 0.82 $\pm$ 0.02  & 0.87 $\pm$ 0.01& \bf 0.90 $\pm$ 0.003 \\
\midrule
rFF + Pooling& 0.7951 $\pm$ 0.0166  & 0.8551 $\pm$ 0.0142 & 0.8933 $\pm$ 0.0156  \\
\midrule
rFF + PMA (ours) & 0.8076 $\pm$ 0.0160 & 0.8534 $\pm$ 0.0152  & 0.8628 $\pm$ 0.0136\\
ISAB (16) + Pooling (ours) & 0.8273 $\pm$ 0.0159 & \bf 0.8915 $\pm$ 0.0144 & \bf 0.9040 $\pm$ 0.0173  \\
ISAB (16) + PMA (ours) & \bf 0.8454 $\pm$ 0.0144 & 0.8662 $\pm$ 0.0149 & 0.8779 $\pm$ 0.0122  \\
\bottomrule
\end{tabular}
\label{table:pointcloud}
\end{table*}

We report the performance of the oracle along with the set-input neural networks in Table~\ref{tab:meta_clustering}.
We additionally report scores of all models after a single EM update.
Overall, the Set Transformer found accurate parameters and even outperformed the oracles after a single EM update.
This may be due to the relatively small size of the input sets; some clusters have fewer than 10 points.
In this regime, sample statistics can differ substantially from population statistics,
which limits the performance of the oracle while the Set Transformer can adapt accordingly.
Notably, the Set Transformer with only 16 inducing points showed the best performance, even outperforming the full Set Transformer.
We believe this is due to the knowledge transfer and regularization via inducing points, helping the network to learn global structures.
Our results also imply that the improvement from using the PMA is more significant than that of the SAB, supporting our claim of the importance of attention-based decoders.
We provide detailed generative processes, network architectures, and training schemes along with additional experiments with various numbers of inducing points in the supplementary material.

\subsection{Set Anomaly Detection}
\label{subsec:meta_anomaly}

\begin{figure}[t]
\centering
\includegraphics[width=0.90\linewidth]{figs/celeba.pdf}
\caption{
Sampled datasets.
Each row is a dataset, consisting of 7 normal images and 1 anomaly (red box).
In each subsampled dataset, a normal image has two attributes (rightmost column) which anomalies do not.
}
\label{fig:celeba_figures}
\end{figure}

\begin{table}[t]
\centering
\small
\caption{Meta set anomaly results.
Each architecture is evaluated using average of test AUROC and test AUPR.
%Random guess shows AUROC as 0.5 and AUPR as 0.125.
%Other describing methods are same with the above description.
%200 test datasets are used to measure the performance of each model and all experiments are repeated 5 times.
}
\vspace{5pt}
\begin{tabular}{@{}ccccc@{}}\toprule
Architecture & Test AUROC & Test AUPR  \\
\midrule
Random guess & 0.5 & 0.125 \\
rFF + Pooling & 0.5643 $\pm$ 0.0139 & 0.4126 $\pm$ 0.0108 \\
rFFp-mean + Pooling & 0.5687 $\pm$ 0.0061 & 0.4125 $\pm$ 0.0127 \\
rFFp-max + Pooling & 0.5717 $\pm$ 0.0117 & 0.4135 $\pm$ 0.0162 \\
rFF + Dotprod & 0.5671 $\pm$ 0.0139 & 0.4155 $\pm$ 0.0115 \\
\midrule
SAB + Pooling (ours) & 0.5757 $\pm$ 0.0143 & 0.4189 $\pm$ 0.0167 \\
%ISAB (4) + Pooling & 0.5679 $\pm$ 0.0155 & 0.4202 $\pm$ 0.0166 \\
rFF + PMA (ours) & 0.5756 $\pm$ 0.0130 & 0.4227 $\pm$ 0.0127 \\
SAB + PMA (ours) & \bf 0.5941 $\pm$ 0.0170 & \bf 0.4386 $\pm$ 0.0089 \\
%Set Transformer (4) & 0.5629 $\pm$ 0.0167 & 0.4122 $\pm$ 0.0151 \\
\bottomrule
\end{tabular}
\label{tab:meta_anomaly}
\end{table}
We evaluate our methods on the task of meta-anomaly detection within a set using the CelebA dataset.
The dataset consists of 202,599 images with the total of 40 attributes.
We randomly sample 1,000 sets of images.
For every set, we select two attributes at random and construct the set by selecting seven images containing both attributes and one image with neither.
The goal of this task is to find the image that does not belong to the set.
We give a detailed description of the experimental setup in the supplementary material.
We report the area under receiver operating characteristic curve (AUROC) and area under precision-recall curve (AUPR) in Table~\ref{tab:meta_anomaly}.
Set Transformers outperformed all other methods by a significant margin.

\subsection{Point Cloud Classification}

We evaluated Set Transformers on a classification task using the ModelNet40~\citep{Chang2015} dataset\footnote{The point-cloud dataset used in this experiment was obtained directly from the authors of \citet{Zaheer2017}.}, which contains three-dimensional objects in $40$ different categories.
Each object is represented as a point cloud, which we treat as a set of $n$ vectors in $\R^3$.
We performed experiments with input sets of size $n \in \{100, 1000, 5000\}$.
Because of the large set sizes, MABs are prohibitively time-consuming due to their $\calO(n^2)$ time complexity.

Table~\ref{table:pointcloud} shows classification accuracies.
We point out that \citet{Zaheer2017} used significantly more engineering for the $5000$ point experiment.
For this experiment only, they augmented data (scaling, rotation) and used a different optimizer (Adamax) and learning rate schedule.
Set Transformers were superior when given small sets, but were outperformed by ISAB (16) + Pooling on larger sets.
First note that classification is harder when given fewer points.
We think Set Transformers were outperformed in the problems with large sets because such sets already had sufficient information for classification, diminishing the need to model complex interactions among points.
We point out that PMA outperformed simple pooling in all other experiments.
